\begin{topicbox}{The Basic Pieces of the Line}
\begin{exposition}
Intervals are the natural named pieces of a number line. This section gives
each basic interval shape its notation and spells out which endpoints are
included.
\end{exposition}
\end{topicbox}
\begin{definitionbox}{Definition (Interval)}
\begin{definition}[Interval]
\label{def:interval}
A set $I\subseteq\mathbb{R}$ is an \emph{interval} if whenever $x,z\in I$ and
$x<y<z$, then $y\in I$.
\end{definition}
\end{definitionbox}
\begin{remark*}[Standard quantified statement]
\[
I\text{ is an interval}
\Longleftrightarrow
\forall x,z\in I\;\forall y\in\mathbb{R}\;
\bigl(x<y<z\Rightarrow y\in I\bigr).
\]
\end{remark*}
\begin{remark*}[Predicate reading]
\[
\operatorname{Interval}(I)
\Longleftrightarrow
\forall x,z\in I\;\forall y\in\mathbb{R}\;
\operatorname{Between}(x,y,z)\Rightarrow y\in I.
\]
\end{remark*}
\begin{remark*}[Negated quantified statement]
\[
I\text{ is not an interval}
\Longleftrightarrow
\exists x,z\in I\;\exists y\in\mathbb{R}\;
\bigl(x<y<z\land y\notin I\bigr).
\]
\end{remark*}
\begin{remark*}[Negation predicate reading]
\[
\neg\operatorname{Interval}(I)
\Longleftrightarrow
\exists x,z\in I\;\exists y\in\mathbb{R}\;
\bigl(\operatorname{Between}(x,y,z)\land y\notin I\bigr).
\]
\end{remark*}
\begin{remark*}[Interpretation]
An interval has no order gaps: once it contains two endpoints, it contains
every real point between them.
\end{remark*}
\begin{dependencies}
\hyperref[def:reals]{Real Numbers};
\hyperref[def:strict-order-on-r]{Strict Order on $\mathbb{R}$}.
\end{dependencies}

\begin{definitionbox}{Definition (Open Interval)}
\begin{definition}[Open Interval]
\label{def:open-interval}
Let $a,b\in\mathbb{R}$ with $a\le b$. The \emph{open interval} with endpoints $a$ and $b$ is
\[
(a,b)=\{\,x\in\mathbb{R} : a<x<b\,\}.
\]
\end{definition}
\end{definitionbox}
\begin{remark*}[Standard quantified statement]
\[
(a,b)=\{\,x\in\mathbb{R} : a<x<b\,\}.
\]
\end{remark*}
\begin{remark*}[Interpretation]
The open interval excludes both endpoints. It is one of the four bounded interval shapes.
\end{remark*}
\begin{dependencies}
\NoLocalDependencies
\end{dependencies}
\begin{definitionbox}{Definition (Closed Interval)}
\begin{definition}[Closed Interval]
\label{def:closed-interval}
Let $a,b\in\mathbb{R}$ with $a\le b$. The \emph{closed interval} with endpoints $a$ and $b$ is
\[
[a,b]=\{\,x\in\mathbb{R} : a\le x\le b\,\}.
\]
\end{definition}
\end{definitionbox}
\begin{remark*}[Standard quantified statement]
\[
[a,b]=\{\,x\in\mathbb{R} : a\le x\le b\,\}.
\]
\end{remark*}
\begin{remark*}[Interpretation]
The closed interval includes both endpoints. It is one of the four bounded interval shapes.
\end{remark*}
\begin{dependencies}
\NoLocalDependencies
\end{dependencies}
\begin{definitionbox}{Definition (Left-Half-Open Interval)}
\begin{definition}[Left-Half-Open Interval]
\label{def:left-half-open-interval}
Let $a,b\in\mathbb{R}$ with $a\le b$. The \emph{left-half-open interval} with endpoints $a$ and $b$ is
\[
(a,b]=\{\,x\in\mathbb{R} : a<x\le b\,\}.
\]
\end{definition}
\end{definitionbox}
\begin{remark*}[Standard quantified statement]
\[
(a,b]=\{\,x\in\mathbb{R} : a<x\le b\,\}.
\]
\end{remark*}
\begin{remark*}[Interpretation]
The left-half-open interval excludes the left endpoint, includes the right. It is one of the four bounded interval shapes.
\end{remark*}
\begin{dependencies}
\NoLocalDependencies
\end{dependencies}

\begin{definitionbox}{Definition (Right-Half-Open Interval)}
\begin{definition}[Right-Half-Open Interval]
\label{def:right-half-open-interval}
Let $a,b\in\mathbb{R}$ with $a\le b$. The \emph{right-half-open interval} with endpoints $a$ and $b$ is
\[
[a,b)=\{\,x\in\mathbb{R} : a\le x<b\,\}.
\]
\end{definition}
\end{definitionbox}
\begin{remark*}[Standard quantified statement]
\[
[a,b)=\{\,x\in\mathbb{R} : a\le x<b\,\}.
\]
\end{remark*}
\begin{remark*}[Interpretation]
The right-half-open interval includes the left endpoint, excludes the right. It is one of the four bounded interval shapes.
\end{remark*}
\begin{dependencies}
\NoLocalDependencies
\end{dependencies}

\begin{theorembox}{Theorem (Characterization of Intervals)}
\begin{theorem}[Characterization of Intervals]
\label{thm:interval-characterization}
A subset $I\subseteq\mathbb{R}$ is an interval if and only if for all $x,z\in I$
with $x<z$,
\[
[x,z]\subseteq I.
\]
\hyperref[prf:interval-characterization]{Go to proof.}
\end{theorem}
\end{theorembox}
\begin{remark*}[Standard quantified statement]
\[
\operatorname{Interval}(I)
\Longleftrightarrow
\forall x,z\in I\;\bigl(x<z\Rightarrow [x,z]\subseteq I\bigr).
\]
\end{remark*}
\begin{remark*}[Predicate reading]
\[
\operatorname{Interval}(I)
\Longleftrightarrow
\operatorname{ClosedSegmentConvex}(I).
\]
\end{remark*}
\begin{remark*}[Negated quantified statement]
\[
\neg\operatorname{Interval}(I)
\Longleftrightarrow
\exists x,z\in I\;\bigl(x<z\land [x,z]\nsubseteq I\bigr).
\]
\end{remark*}
\begin{remark*}[Negation predicate reading]
\[
\exists x,z\in I\;\bigl(x<z\land [x,z]\nsubseteq I\bigr)
\Longleftrightarrow
\neg\operatorname{Interval}(I).
\]
\end{remark*}
\begin{remark*}[Interpretation]
The pointwise betweenness definition of interval is equivalent to saying that
every closed segment joining two points of the set stays inside the set.
\end{remark*}
\begin{dependencies}
\hyperref[def:interval]{Interval};
\hyperref[def:closed-interval]{Closed Interval}.
\end{dependencies}

\begin{theorembox}{Theorem (Classification of Intervals in $\mathbb{R}$)}
\begin{theorem}[Classification of Intervals in $\mathbb{R}$]
\label{thm:classification-of-intervals-in-r}
Every interval in $\mathbb{R}$ is one of the standard interval types:
\[
(a,b),\ [a,b],\ (a,b],\ [a,b),\ (a,\infty),\ [a,\infty),
\]
\[
(-\infty,b),\ (-\infty,b],\ \mathbb{R},\ \varnothing,
\]
including degenerate one-point intervals.
\hyperref[prf:classification-of-intervals-in-r]{Go to proof.}
\end{theorem}
\end{theorembox}
\begin{remark*}[Standard quantified statement]
\[
\forall I\subseteq\mathbb{R}\;
\bigl(\operatorname{Interval}(I)\Rightarrow
I\text{ is one of the standard interval types}\bigr).
\]
\end{remark*}
\begin{remark*}[Predicate reading]
\[
\operatorname{Interval}(I)\Longrightarrow \operatorname{StandardRealInterval}(I).
\]
\end{remark*}
\begin{remark*}[Negated quantified statement]
\[
\exists I\subseteq\mathbb{R}\;
\bigl(\operatorname{Interval}(I)\land
\neg\operatorname{StandardRealInterval}(I)\bigr).
\]
\end{remark*}
\begin{remark*}[Negation predicate reading]
\[
\neg\operatorname{StandardRealInterval}(I)\Longrightarrow \neg\operatorname{Interval}(I).
\]
\end{remark*}
\begin{remark*}[Interpretation]
The real-line interval condition has no exotic bounded or unbounded shapes:
the endpoint data determine one of the familiar open, closed, half-open, ray,
whole-line, empty, or singleton forms.
\end{remark*}
\begin{dependencies}
\hyperref[def:interval]{Interval};
\hyperref[def:open-interval]{Open Interval};
\hyperref[def:closed-interval]{Closed Interval};
\hyperref[def:left-half-open-interval]{Left-Half-Open Interval};
\hyperref[def:right-half-open-interval]{Right-Half-Open Interval};
\hyperref[def:ray]{Ray}.
\end{dependencies}

\begin{figure}[H]
\centering
\begin{tikzpicture}[x=1.0cm,y=0.8cm,>=Latex]
  \draw[->] (-0.4,0) -- (7.6,0) node[right] {$\mathbb R$};
  \foreach \x/\lab in {1.5/$a$,5.4/$b$} {
    \draw (\x,0.1) -- (\x,-0.1) node[below=3pt] {\lab};
  }
  \draw[line width=1.2pt] (1.5,0.9) -- (5.4,0.9);
  \fill[white,draw=black] (1.5,0.9) circle (2pt);
  \fill[white,draw=black] (5.4,0.9) circle (2pt);
  \node[right] at (5.55,0.9) {$(a,b)$};
  \draw[line width=1.2pt] (1.5,1.6) -- (5.4,1.6);
  \fill (1.5,1.6) circle (2pt);
  \fill (5.4,1.6) circle (2pt);
  \node[right] at (5.55,1.6) {$[a,b]$};
  \draw[line width=1.2pt] (1.5,2.3) -- (7.1,2.3);
  \fill (1.5,2.3) circle (2pt);
  \draw[->,line width=1.2pt] (5.4,2.3) -- (7.1,2.3);
  \node[right] at (7.15,2.3) {$[a,\infty)$};
  \draw[->,line width=1.2pt] (1.5,3.0) -- (-0.1,3.0);
  \fill[white,draw=black] (5.4,3.0) circle (2pt);
  \draw[line width=1.2pt] (-0.1,3.0) -- (5.4,3.0);
  \node[right] at (5.55,3.0) {$(-\infty,b)$};
\end{tikzpicture}
\caption{The classification theorem says that intervals in $\mathbb R$ are built from endpoint inclusion choices and possible unbounded directions.}
\label{fig:interval-classification}
\end{figure}

\begin{definitionbox}{Definition (Ray)}
\begin{definition}[Ray]
\label{def:ray}
Let $a\in\mathbb{R}$. The four \emph{rays} determined by $a$ are
\[
(a,\infty)=\{x:x>a\},\ [a,\infty)=\{x:x\ge a\},\ (-\infty,a)=\{x:x<a\},\ (-\infty,a]=\{x:x\le a\}.
\]
\end{definition}
\end{definitionbox}
\begin{remark*}[Standard quantified statement]
\[
\begin{gathered}
(a,\infty)=\{\,x\in\mathbb{R}:a<x\,\},\qquad
[a,\infty)=\{\,x\in\mathbb{R}:a\le x\,\},\\
(-\infty,a)=\{\,x\in\mathbb{R}:x<a\,\},\qquad
(-\infty,a]=\{\,x\in\mathbb{R}:x\le a\,\}.
\end{gathered}
\]
\end{remark*}
\begin{remark*}[Interpretation]
Rays are the unbounded one-sided intervals. Together with the bounded intervals and the whole line $(-\infty,\infty)=\mathbb{R}$, they exhaust the interval shapes.
\end{remark*}
\begin{dependencies}
\NoLocalDependencies
\end{dependencies}
